\begin{helper}
Let \(V\) be finite and let \(\nu\) be a measure on activation--priority
pairs \((A,\pi)\in\mathcal P(V)\times\mathbb R^V\).  Fix \(0\le p\le1\),
distinct \(u,v\in V\), and disjoint finite sets
\(U,T\subseteq V\setminus\{u,v\}\).  Assume the activation atom law
\[
\nu(A=A_0)=\operatorname{ofReal}\bigl(p^{|A_0|}(1-p)^{|V|-|A_0|}\bigr)
\]
for every \(A_0\subseteq V\).  Assume also that, on each activation atom, the
activated priorities are supported in \([0,1]\), and that every full lower
rectangle satisfies
\[
\nu(A=A_0,\ 0\le \pi(q)\le t_q\text{ for all }q\in V)
=\nu(A=A_0)\operatorname{ofReal}\!\left(\prod_{q\in V}t_q\right)
\]
whenever \(0\le t_q\le1\) for all \(q\).

Then the ordered chamber
\[
\begin{aligned}
E_{u,v}(U,T)=\{(A,\pi):\;&u,v\in A,\ \pi(v)\le\pi(u),\
\pi(q)\in[0,1]\text{ for every }q\in A,\\
& z\in U\cap A\Rightarrow \pi(z)<\pi(u),\
z\in T\cap A\Rightarrow \pi(z)<\pi(v)\}
\end{aligned}
\]
has mass
\[
\nu(E_{u,v}(U,T))
=\operatorname{ofReal}\!\left(
\int_0^1\int_0^a
p^2\bigl((1-p)+pa\bigr)^{|U|}
\bigl((1-p)+pb\bigr)^{|T|}\,db\,da\right).
\]
\end{helper}
\begin{proof}
Put \(S=U\cup T\) and \(C=S\cup\{u,v\}\).  Since \(u\ne v\),
\(u,v\notin S\), and \(U\cap T=\emptyset\), the sets
\(U,T,\{u\},\{v\}\) are pairwise disjoint.  Fix an activation atom
\(A_0\).  Applying the lower-rectangle hypothesis with all thresholds equal
to \(1\) shows that the full cube
\(\{0\le \pi(q)\le1\text{ for every }q\in V\}\) has the whole mass of the
atom \(A=A_0\).  Thus restricting an atom to the full cube changes no
\(\nu\)-measure.

On this full cube, the lower-rectangle hypothesis gives the masses of all
anchored rectangles \(\prod_q[0,t_q]\), \(0\le t_q\le1\), as the activation
atom mass times \(\prod_q t_q\).  These anchored rectangles form a
\(\pi\)-system generating the Borel sigma-algebra of the finite cube
\([0,1]^V\), and both finite measures have the same total mass.  The
finite-measure uniqueness theorem therefore identifies the priority law on
the atom with the activation atom mass times finite product Lebesgue measure
on \([0,1]^V\).  This is the same lower-rectangle-to-cell calculation
registered in \noderef{SamplingFiniteCoordinateMarginalCellLaw}, with the
atomwise interval calculation supplied by
\noderef{SamplingFinitePriorityOneCoordinateIntervalMass}.

Partition the chamber by the activated subsets \(B_U\subseteq U\),
\(B_T\subseteq T\), and \(R\subseteq V\setminus C\).  The corresponding atom
is
\[
A_{B_U,B_T,R}=\{u,v\}\cup B_U\cup B_T\cup R .
\]
The associated chamber pieces are pairwise disjoint, and their union differs
from \(E_{u,v}(U,T)\) only by atomwise complements of the activated-priority
cube, which have measure \(0\).  Hence finite additivity reduces
\(\nu(E_{u,v}(U,T))\) to the finite sum of these piece masses.

For a fixed triple \((B_U,B_T,R)\), write \(a=\pi(u)\) and \(b=\pi(v)\).
The product-measure identification on the atom and Tonelli's theorem on the
finite cube give the piece volume
\[
\int_0^1\int_0^a a^{|B_U|}b^{|B_T|}\,db\,da .
\]
Indeed, once \(0\le b\le a\le1\) is fixed, each coordinate in \(B_U\) ranges
independently over an interval of length \(a\), each coordinate in \(B_T\)
ranges independently over an interval of length \(b\), and all remaining
priority coordinates have interval length \(1\).  Thus the piece mass is
\[
\operatorname{ofReal}\!\left(
p^{|A_{B_U,B_T,R}|}(1-p)^{|V|-|A_{B_U,B_T,R}|}
\int_0^1\int_0^a a^{|B_U|}b^{|B_T|}\,db\,da\right).
\]

Now sum first over \(R\subseteq V\setminus C\).  Since the displayed sets are
disjoint from \(V\setminus C\),
\[
\sum_R p^{|R|}(1-p)^{|V\setminus C|-|R|}
=(p+(1-p))^{|V\setminus C|}=1 .
\]
The remaining finite sums over \(B_U\subseteq U\) and \(B_T\subseteq T\) are
nonnegative, so finite sums may be moved through the two integrals.  The
\(U\)-sum is the binomial expansion
\[
\sum_{B_U\subseteq U}p^{|B_U|}(1-p)^{|U|-|B_U|}a^{|B_U|}
=((1-p)+pa)^{|U|},
\]
and the \(T\)-sum is
\[
\sum_{B_T\subseteq T}p^{|B_T|}(1-p)^{|T|-|B_T|}b^{|B_T|}
=((1-p)+pb)^{|T|}.
\]
Multiplying by the forced activations of \(u\) and \(v\) gives the factor
\(p^2\).  Substituting these two binomial identities into the finite additive
sum gives exactly the stated integral formula.
\end{proof}
